\documentclass[11pt]{article}
\usepackage{amsmath,amssymb,amsfonts,amsthm,booktabs,array,tabularx,geometry,hyperref}
\geometry{margin=2.7cm}
\hypersetup{colorlinks=true,linkcolor=blue,citecolor=blue,urlcolor=blue}

\newtheorem{theorem}{Theorem}
\newtheorem{remark}{Remark}
\newcommand{\R}{\mathbb R}
\newcommand{\pos}[1]{\left[#1\right]_+}

\title{Componentwise Sustainability Accounting:\\
Typed Stock--Flow Ledgers and Depletion Diagnostics}
\author{Anonymous}
\date{}

\begin{document}
\maketitle

\begin{abstract}
Sustainability assessments often track heterogeneous physical stocks and service flows but summarize them with a compensatory scalar index. We develop a deliberately limited accounting framework that separates three questions: whether current service supply meets current demand, whether each physical compartment obeys an admissible stock--flow balance, and whether a specified trajectory remains inside declared safety constraints. First, we record the elementary, domain-qualified observation that a positively weighted scalar balance cannot certify every componentwise balance on an unrestricted balance space. Second, we formulate a typed incidence/stoichiometric ledger scaffold and instantiate it as a six-compartment donor-limited example for one conserved material. A typed service readout links selected physical fluxes to contemporaneous balances, while a service-set comparison distinguishes current delivery from feasibility under declared regenerative, boundary, quality, and exergy restrictions; neither establishes net depletion. Conservation and forward invariance of the example are established separately. Third, we distinguish gross turnover, a local net-depletion ratio, and a scenario-conditioned hitting time. Positive gross uptake establishes throughput or dependency, not net depletion. We illustrate the terminology with a groundwater linear-trend anomaly index, a phosphate reserve-life ratio, and a separately classified removals-only fisheries pressure time scale. The framework is one accounting layer within a wider family of sustainability models; it is not a completed general theory, a demographic forecasting model, or a derivation of any particular delayed-feedback core.
\end{abstract}

\section{Scope and contribution}

Vector accounting is already standard in material-flow analysis, environmental input--output analysis, life-cycle assessment, and multi-capital accounting \cite{leontief1970,iso14040,seea2014,dasgupta2021}. Official economy-wide material-flow guidance likewise organizes physical flows by material category and system boundary \cite{eurostat2001}. The problem addressed here is narrower than replacing those traditions. We ask what can and cannot be certified from componentwise physical accounts when the components have different units, thresholds, and replenishment mechanisms.

The paper makes three contributions.

\begin{enumerate}
    \item It states an elementary, domain-qualified non-compensation observation. A positive scalar aggregate can summarize a vector, but it cannot by itself certify that every component is non-negative on an unrestricted balance space.
    \item It gives a typed incidence/stoichiometric scaffold and a six-compartment donor-limited illustration in which conservation and physical non-negativity are distinct properties. An explicit service readout shows how selected ledger fluxes can generate contemporaneous balance components without identifying internal ecosystem transfers with services. A support-provenance interface contrasts the all-pathway service set with the subset satisfying declared regeneration, boundary, material-grade, and exergy or capacity constraints.
    \item It separates gross dependency, local net decline, and dynamical time to a threshold. These are useful but different quantities and should not share one depletion-horizon label; an isolated gross-removal time scale is classified outside that net-depletion hierarchy.
\end{enumerate}

This accounting framework is one component of a larger model family. Delayed effort dynamics, bifurcation analysis, robust viability, stage structure, spatial moments, and thermodynamic constraints require separate dynamical or theoretical treatment. No reduction from this ledger to those models is asserted.

\section{Typed states, services, and contemporaneous balances}

\subsection{Physical states and service observations}

Let
\[
x(t)=(x_1(t),\ldots,x_m(t))\in X\subseteq\R_+^m
\]
be a vector of typed physical stocks. A type records at least the material or ecological identity, compartment, spatial support, and physical unit. Quantities are added in a conservation equation only when they represent the same conserved moiety or have been converted by an explicit stoichiometric map.

Services are observations or feasible outputs of the physical state rather than additional conserved mass. For services indexed by $i=1,\ldots,n$, write
\[
s_i(t)=\mathcal S_{\rm svc}_i(x(t),u(t),\theta),
\]
where $u$ denotes admissible operating or extraction choices and $s_i$ and the corresponding demand $d_i$ have the same service-specific units. The observation operator $\mathcal S_{\rm svc}$ may be linear in a restricted application, but linearity is not required by the accounting logic. Internal physical transfers are not services merely because they appear in a ledger: a typed readout must identify the delivered flow, its boundary, and any unit conversion.

Define the contemporaneous component balance
\begin{equation}
 b_i(t)=s_i(t)-d_i(t).
 \label{eq:component-balance}
\end{equation}
Then $b_i(t)\ge0$ means that measured supply meets measured demand for component $i$ at that instant. It does \emph{not} by itself imply that the underlying trajectory is sustainable. A stock can meet current demand while declining toward a threshold, and a stock below a desired level can have a positive current flow balance while recovering.

For an admissible operating set $\mathcal U(x,t)$ and a declared demand set $\mathcal D(t)$, the physically attainable balance domain is
\begin{equation}
\mathcal B(x,t)
=\{\mathcal S_{\rm svc}(x,u,\theta)-d:
      u\in\mathcal U(x,t),\ d\in\mathcal D(t)\}.
\label{eq:feasible-balance-domain}
\end{equation}
Thus the geometry of the balance domain is state dependent and inherited partly from the stock--flow model. The unrestricted argument below cannot replace an application-specific analysis of $\mathcal B(x,t)$.

\subsection{Service adequacy and support provenance}

Current service adequacy and regenerative service feasibility are different claims. Let $\Gamma_{\rm all}(x,t)\subseteq\R_+^n$ contain the service vectors feasible through all pathways admitted by an application, and let
\[
\Gamma_{\rm reg}(x,t)\subseteq\Gamma_{\rm all}(x,t)
\]
be the feasible set after imposing the declared regenerative-flow, system-boundary, material-quality, and exergy or capacity restrictions. These correspondences are application inputs obtained from a typed pathway or technology model; the stock ledger alone does not construct them.

Assume $0\in\Gamma_{\rm reg}(x,t)$ and choose a nonzero service direction $\bar s\in\Gamma_{\rm all}(x,t)$. Define
\begin{equation}
\alpha_{\rm reg}(\bar s;x,t)
=\sup\{\alpha\in[0,1]:\alpha\bar s\in\Gamma_{\rm reg}(x,t)\}.
\label{eq:regenerative-support}
\end{equation}
If $\Gamma_{\rm reg}(x,t)$ is closed along the ray $\{\alpha\bar s:0\le\alpha\le1\}$, the supremum is attained. Closedness alone does not imply that every smaller fraction is feasible. That radial interval property follows, for example, if $\Gamma_{\rm reg}(x,t)$ is star-shaped about the origin (convexity with $0\in\Gamma_{\rm reg}$ is sufficient). Without these assumptions, $\alpha_{\rm reg}$ is only a supremal directional support fraction.
The vector $(1-\alpha_{\rm reg})\bar s$ is a directional support gap measured in the same service units as $\bar s$. A realized service $s\in\Gamma_{\rm all}\setminus\Gamma_{\rm reg}$ is support-dependent relative to the declared pathway restrictions even when $s\ge d$. This statement neither subtracts raw material from service nor proves that a physical stock is declining; net depletion still requires a negative stock balance or a trajectory argument.

\subsection{What the scalar aggregate cannot certify}

Compensability is a long-standing design issue in composite indicators, and non-compensatory ranking methods already exist \cite{munda2009}. The observation below makes a narrower logical claim: unless the feasible balance domain supplies additional restrictions, a positive weighted sum cannot certify the conjunction of componentwise non-negativity. This certification question is related to, but does not by itself settle, the debate between weak and strong sustainability or identify which stocks constitute critical natural capital \cite{neumayer2013,ekins2003}.

A common summary is
\[
\sigma_w(b)=w^\top b,
\qquad w\in\R_{++}^n.
\]
The relevant geometry is elementary. Let $n\ge2$, choose a component $k$, and prescribe any deficit $L>0$. For any $j\ne k$, the vector
\[
b_k=-L,
\qquad
b_j=\frac{w_kL+1}{w_j},
\qquad
b_i=0\quad(i\notin\{j,k\})
\]
has $w^\top b=1>0$ despite $b_k<0$. A positive half-space therefore contains vectors outside the non-negative orthant. This familiar fact is recorded here only to fix the certification boundary, not as a new mathematical result.

On a restricted feasible domain such as Eq.~\eqref{eq:feasible-balance-domain}, a scalar certificate would require the separately established implication
\[
b\in\mathcal B(x,t),
\quad w^\top b\ge0
\quad\Longrightarrow\quad
b\in\R_+^n.
\]
The implication fails whenever the feasible domain contains a compensating counterexample of the displayed kind. If physical restrictions exclude all such vectors, the unrestricted construction decides nothing: the implication must be proved from those restrictions. This domain qualification is essential because arbitrarily large compensating surpluses need not be physically attainable in a particular application.

The observation rules out an \emph{unrestricted compensatory certificate}; it does not imply that every scalar statistic is uninformative. In particular,
\[
\min_i b_i,
\qquad
\|[-b]_+\|,
\qquad\text{and}\qquad
\max_i\frac{[-b_i]_+}{s_i^{\rm ref}}
\]
are non-compensatory scalar encodings when the positive reference scales $s_i^{\rm ref}$ and zero conventions are declared. They can certify absence of a component deficit because they preserve the conjunction $b_i\ge0$ rather than trading one component against another. They still discard the identity and cause of the limiting component, so the vector should accompany them. Scalar summaries may rank or communicate; certification requires the full vector, a logically equivalent non-compensatory encoding, component thresholds, or a proved feasible-domain implication.

\section{A typed flux-ledger scaffold and a six-compartment example}

The preceding observation concerns certification of contemporaneous service balances. We now turn to a physical stock--flow ledger, whose central claims concern conservation and boundary admissibility. The abstract scaffold and the concrete illustration are distinguished explicitly.

\subsection{Incidence, types, and service readouts}

Let $z\in\R_+^m$ collect compartments after all entries have been assigned a material identity, spatial support, and unit. Let $v(z,u)\in\R_+^r$ be non-negative primitive fluxes and let $S\in\R^{m\times r}$ record their signed incidence or stoichiometric coefficients:
\begin{equation}
\dot z=S v(z,u).
\label{eq:typed-ledger}
\end{equation}
Entries may be added within a row only when their types and units agree; a conversion between types must be represented by an explicit stoichiometric coefficient rather than an implicit sum. For one closed conserved moiety with conversion vector $\ell\in\R_{++}^m$, conservation is the left-nullspace condition $\ell^\top S=0$. For several conserved moieties, a conservation matrix $L$ replaces $\ell^\top$. This is the standard stoichiometric-subspace bookkeeping of reaction-network theory \cite{feinberg2019}; no mass-action assumption is imposed on $v$.

For flux $j$, let $D_j=\{i:S_{ij}<0\}$ denote its donor coordinates. The donor-face condition
\[
z_i=0\text{ for some }i\in D_j
\quad\Longrightarrow\quad
v_j(z,u)=0
\]
is a transparent sufficient condition for non-negative-cone invariance: on the face $z_i=0$, every term that would decrease coordinate $i$ vanishes. Conservation and this boundary condition are separate obligations.

A typed service readout is an additional map, not part of mass conservation. In applications where delivered services are selected or converted ledger fluxes, one may write
\begin{equation}
s=\mathcal S_{\rm svc}(z,u,\theta)=Q(\theta)v(z,u),
\label{eq:service-readout}
\end{equation}
where every row of $Q$ declares the delivery boundary and conversion into one service-specific unit. More general state-dependent readouts are possible. Combining Eqs.~\eqref{eq:typed-ledger} and \eqref{eq:service-readout} generates both physical balances and the attainable service-balance domain in Eq.~\eqref{eq:feasible-balance-domain}; it does not yet establish trajectory safety.

\subsection{Six-compartment illustration}

For one conserved limiting material, consider the compartments

\begin{center}
\begin{tabular}{cl}
$X$ & living biomass,\\
$U$ & detritus or recoverable residual,\\
$A$ & active abiotic pool,\\
$G$ & geological or slowly available pool,\\
$P$ & product or in-use stock,\\
$W$ & absorbing or currently unavailable stock.
\end{tabular}
\end{center}

All primitive fluxes are non-negative on the admissible state domain:

\begin{center}
\begin{tabular}{cl}
$g(X,A)$ & assimilation from $A$ to $X$,\\
$m(X)$ & mortality/turnover from $X$ to $U$,\\
$h(X,E)$ & harvest from $X$,\\
$d_U(U)$ & decomposition from $U$ to $A$,\\
$e_{GA}(G,A)$ & geological-to-active transfer,\\
$e_{AG}(A,G)$ & active-to-geological transfer,\\
$c_G(G,E_G)$ & direct mining from $G$,\\
$r_P(P)$ & product retirement.
\end{tabular}
\end{center}

Let $\alpha\in[0,1]$ be the fraction of harvest entering $U$ immediately, and $\rho\in[0,1]$ the fraction of retired product that returns to $U$ rather than $W$. A closed six-compartment illustration is

\begin{align}
\dot X &= g(X,A)-m(X)-h(X,E), \label{eq:X}\\
\dot U &= m(X)+\alpha h(X,E)+\rho r_P(P)-d_U(U), \label{eq:U}\\
\dot A &= -g(X,A)+d_U(U)+e_{GA}(G,A)-e_{AG}(A,G), \label{eq:A}\\
\dot G &= -e_{GA}(G,A)+e_{AG}(A,G)-c_G(G,E_G), \label{eq:G}\\
\dot P &= (1-\alpha)h(X,E)+c_G(G,E_G)-r_P(P), \label{eq:P}\\
\dot W &= (1-\rho)r_P(P). \label{eq:W}
\end{align}

With
\[
z=(X,U,A,G,P,W)^\top,
\quad
v=(g,m,h,d_U,e_{GA},e_{AG},c_G,r_P)^\top,
\]
these equations have the form \eqref{eq:typed-ledger} with
\[
S(\alpha,\rho)=
\begin{pmatrix}
 1&-1&-1& 0& 0& 0& 0& 0\\
 0& 1&\alpha&-1& 0& 0& 0&\rho\\
-1& 0& 0& 1& 1&-1& 0& 0\\
 0& 0& 0& 0&-1& 1&-1& 0\\
 0& 0&1-\alpha&0&0&0&1&-1\\
 0& 0& 0& 0& 0& 0& 0&1-\rho
\end{pmatrix},
\qquad \mathbf 1^\top S=0.
\]
The matrix makes the routing choices visible: the constant splits $\alpha$ and $\rho$, the compartment set, and the absorbing-sink convention are constitutive features of this example, not properties of every typed ledger. The retirement map is general donor-limited notation; $r_P(P)=\lambda_P P$ is one optional constitutive choice.

This is a monomaterial projection, not a universal ecological mechanism. Coupled carbon, nitrogen, phosphorus, water, or other accounts require additional typed rows and a conservation matrix. The sink $W$ is absorbing in this illustration but, under the conservation and invariance conditions proved below, it is bounded above by $M(0)$. Non-material effects of accumulated sink mass can enter rates such as $g(X,A,W)$ or $m(X,W)$; material remobilization from $W$ requires an additional donor-limited column in $S$.

The labels ``recoverable residual'' and retirement-to-$U$ specify mass routing only. They do not assert perfect recovery, grade-preserving recycling, retained accessibility, or energetic feasibility. If an application makes such claims, $U$ and $P$ must be split by relevant material, location, and quality grade; recovery and upgrading must carry declared yields, residual routes, and exergy or capacity inputs. The conservation argument below then applies to the expanded typed incidence system, not automatically to an undifferentiated quality-neutral loop.

For a minimal service readout, suppose the two declared services are gross harvested material delivered to products and gross mined material delivered to products. Then
\begin{equation}
s_{\rm harv}=(1-\alpha)h(X,E),
\qquad
s_{\rm mine}=c_G(G,E_G),
\label{eq:example-service-readout}
\end{equation}
and $b_{\rm harv}=s_{\rm harv}-d_{\rm harv}$ and $b_{\rm mine}=s_{\rm mine}-d_{\rm mine}$ for demands in matching units. This is an illustrative readout, not a claim that all harvest or mining is final welfare delivery. It makes the coupling explicit: donor limitation forces the corresponding gross delivery to zero at $X=0$ or $G=0$, while internal assimilation $g$ is not classified as a societal service unless an application supplies a separate service definition.

\subsection{Boundary assumptions}

Let $E(t)$ and $E_G(t)$ be prescribed non-negative, continuous, and locally bounded effort histories. Assume
\begin{align*}
g(0,A)&=0, & g(X,0)&=0,\\
m(0)&=0, & h(0,E)&=0,\\
d_U(0)&=0, & e_{GA}(0,A)&=0,\\
e_{AG}(0,G)&=0, & c_G(0,E_G)&=0,\\
r_P(0)&=0,
\end{align*}
with every primitive flux continuous in effort, locally Lipschitz in the state uniformly on compact effort sets, and non-negative on $\R_+^6$. These are donor conditions: a transfer vanishes when its donor is empty. They also make explicit that the displayed model is a non-autonomous ordinary differential equation when the effort histories are prescribed rather than generated by additional state equations.

A target-relaxation law such as $e_{GA}=\omega(A^{\rm eq}-A)$ does not satisfy the finite-donor condition unless it is also limited by $G$. It may be used only if the source is declared an effectively infinite external reservoir, in which case the six-compartment system is open rather than closed.

\begin{theorem}[Mass conservation]
For Eqs.~\eqref{eq:X}--\eqref{eq:W}, the total tracked material
\[
M=X+U+A+G+P+W
\]
is constant along every classical solution.
\end{theorem}

\begin{proof}
Differentiate $M$ and substitute Eqs.~\eqref{eq:X}--\eqref{eq:W}:
\begin{align*}
\dot M
&=(g-m-h)
 +(m+\alpha h+\rho r_P-d_U)\\
&\quad+(-g+d_U+e_{GA}-e_{AG})
 +(-e_{GA}+e_{AG}-c_G)\\
&\quad+((1-\alpha)h+c_G-r_P)
 +(1-\rho)r_P.
\end{align*}
Here each symbol denotes the same primitive flux wherever it appears. The assimilation terms are $g-g=0$; mortality gives $-m+m=0$; decomposition gives $-d_U+d_U=0$; geological exchange gives $e_{GA}-e_{GA}=0$ and $-e_{AG}+e_{AG}=0$; and mining gives $-c_G+c_G=0$. The remaining harvest and retirement terms satisfy
\[
-h+\alpha h+(1-\alpha)h=0,
\qquad
\rho r_P-r_P+(1-\rho)r_P=0.
\]
Consequently $\dot M(t)=0$ at every time at which the classical solution is defined. Integrating from $0$ to $t$ gives $M(t)-M(0)=0$, which proves conservation.
\end{proof}

\begin{theorem}[Forward invariance of the non-negative cone]
Under the boundary assumptions above, $\R_+^6$ is forward invariant for Eqs.~\eqref{eq:X}--\eqref{eq:W}.
\end{theorem}

\begin{proof}
It suffices to verify the inward-pointing condition on every coordinate face. If $X=0$, the donor assumptions give $g(0,A)=m(0)=h(0,E)=0$, and hence $\dot X=0$. If $U=0$, then $d_U(0)=0$ and
\[
\dot U=m(X)+\alpha h(X,E)+\rho r_P\ge0.
\]
If $A=0$, then $g(X,0)=e_{AG}(0,G)=0$, so
\[
\dot A=d_U(U)+e_{GA}(G,0)\ge0.
\]
If $G=0$, then $e_{GA}(0,A)=c_G(0,E_G)=0$, and therefore $\dot G=e_{AG}(A,0)\ge0$. If $P=0$, its retirement outflow vanishes and
\[
\dot P=(1-\alpha)h(X,E)+c_G(G,E_G)\ge0.
\]
Finally, on $W=0$ one has $\dot W=(1-\rho)r_P\ge0$.

Thus the vector field belongs to the tangent cone of $\R_+^6$ at every boundary point. The primitive fluxes are locally Lipschitz by assumption, so the standard tangent-cone invariance theorem for ordinary differential equations applies \cite{aubin1991}. A solution beginning in $\R_+^6$ therefore cannot leave the cone while it exists.
\end{proof}

\begin{remark}[Open systems]
Imports, exports, atmospheric losses, and cross-boundary transport can be added as typed boundary fluxes. The balance then becomes $\dot M=I_{\partial}-O_{\partial}$ rather than zero. Writing these flows explicitly is preferable to preserving a nominal invariant by allowing an unobserved or finite donor compartment to become negative.
\end{remark}

\section{Three different notions of depletion time}

The ledger supplies the net active-pool derivative needed to distinguish gross throughput from net decline and from a model-conditioned threshold time. Let $A_{\min}$ be a declared threshold for the active abiotic pool and assume $A>A_{\min}$.

\subsection{Gross turnover or dependency}

Assimilation $g(X,A)>0$ establishes that the living stock uses material from $A$. Define the gross turnover intensity
\begin{equation}
J_A^{\rm gross}=\frac{g(X,A)}{A},
\label{eq:gross-intensity}
\end{equation}
and, when $g>0$, the gross support-coverage ratio
\begin{equation}
H_A^{\rm gross}=\frac{A-A_{\min}}{g(X,A)}.
\label{eq:gross-coverage}
\end{equation}
Both have useful interpretations, but neither is a time to depletion. At an interior steady state, $g$ can be positive while decomposition and geological transfer balance it exactly, so that $\dot A=0$.

The implication
\[
g>0\quad\Longrightarrow\quad\dot A<0
\]
is false in general. Gross uptake measures throughput or dependency; net depletion is a balance property.

\subsection{Local net-depletion ratio}

A local constant-rate diagnostic is
\begin{equation}
H_A^{\rm loc}(t)=
\frac{A(t)-A_{\min}}{\pos{-\dot A(t)}},
\label{eq:local-horizon}
\end{equation}
with the extended-real convention $H_A^{\rm loc}=+\infty$ when $\dot A\ge0$. This correctly reports no current net decline at a stationary or replenishing state.

Equation~\eqref{eq:local-horizon} is still not a trajectory forecast. It freezes the current net rate. If the fluxes change with $A$, policy, climate, prices, or other states, the realized threshold time can differ substantially.

\subsection{Scenario-conditioned hitting time}

For a fully specified dynamical model, policy/scenario $\pi$, disturbance history $d$, and initial state $x_0$, define
\begin{equation}
T_A(x_0;\pi,d)
=\inf\{t\ge0:A^{\pi,d}(t;x_0)\le A_{\min}\},
\label{eq:hitting-time}
\end{equation}
with $T_A=+\infty$ if the threshold is never reached. Under parameter, observation, and scenario uncertainty, the appropriate output is a distribution or robust interval of $T_A$, not a single universal date.

The three quantities answer different questions:

\begin{center}
\begin{tabularx}{\textwidth}{lX}
\toprule
Quantity & Question\\
\midrule
$J_A^{\rm gross}$ or $H_A^{\rm gross}$ & How strongly does the system depend on or turn over the pool at the current gross rate?\\
$H_A^{\rm loc}$ & If the current net decline were frozen, what is the local stock-to-rate ratio?\\
$T_A$ & Under a stated model, policy, and disturbance scenario, when is the threshold first reached?\\
\bottomrule
\end{tabularx}
\end{center}

\section{Applications and diagnostic interpretation}

These applications apply or delimit the diagnostic terminology. The groundwater and phosphate quantities are narrowly classified accounting ratios, while the fisheries construction is retained specifically to show why an isolated gross-removal time scale must not be promoted to a net depletion diagnostic. None is a calibrated early-warning system.

\subsection{Groundwater anomalies}

The G3P v1.12 product provides monthly groundwater-storage anomalies relative to a reference period rather than absolute aquifer volumes \cite{guentner2024}. For a basin-mean anomaly series over the reported April 2002--September 2023 product window, define the \textbf{Linear-Trend Anomaly Persistence Index}
\begin{equation}
L_{\rm hist}^{\rm anom}
=\frac{a_{\rm latest}-a_{\rm hist,min}}{[-\widehat{\dot a}]_+}.
\label{eq:anomaly-persistence}
\end{equation}
It is the fitted distance to the series' own historical minimum divided by the fitted decline rate. It is therefore a statistical anomaly index with units of time, not the physical stock ratio $H_A^{\rm loc}$ and not a forecast of aquifer exhaustion. Its value depends on the product window, basin mask, anomaly reference, and linear-trend convention.

\begin{table}[h]
\centering
\caption{G3P v1.12 basin-mean anomaly trends and Linear-Trend Anomaly Persistence Indices over the reported April 2002--September 2023 window.}
\begin{tabular}{lcc}
\toprule
Basin & Basin-mean anomaly trend (cm/yr) & $L_{\rm hist}^{\rm anom}$ (yr)\\
\midrule
Indo--Gangetic & $-49.7$ & $\approx2.7$\\
North China Plain & $-18.6$ & $\approx7.9$\\
Central Valley & $-16.1$ & $\approx9.5$\\
La Mancha & $-3.2$ & $\approx21.4$\\
\bottomrule
\end{tabular}
\end{table}

The entries summarize fitted basin means, not pixel extremes. They rank persistence of anomaly decline within the selected product window but do not identify recharge mechanisms, recoverable storage, or an absolute safety threshold. A physical groundwater $H_A^{\rm loc}$ requires an absolute stock estimate and a net stock derivative---for example, aquifer geometry or saturated thickness together with appropriate storage parameters---rather than an anomaly series alone. No nonlinear hydraulic-response or land-subsidence model is fitted here.

\subsection{Phosphate rock}

USGS reports distinguish reserves---currently economic under prevailing conditions---from broader geological resources \cite{usgsmcs}. At constant current production $C_G$, the reserve-life ratio is
\begin{equation}
T_{\rm reserve}=\frac{G_{\rm reserve}}{C_G}.
\end{equation}
Using approximately $74{,}000$ Mt of world reserves and $240{,}000$ kt/yr of production gives approximately $309$ years. This arithmetic is internally consistent as a reserve-life ratio to zero. It is not a physical exhaustion forecast because reserve classification changes with prices, technology, exploration, and regulation.

A resource-threshold calculation such as
\[
T_{\rm resource,10\%}
=\frac{0.9G_{\rm resource}}{C_G}
\]
answers a different question and must not be placed in the same column without an explicit convention label.

\subsection{Fisheries: a removals-only pressure time scale}

Spawning-stock biomass is adult biological biomass. It is not the active abiotic pool $A$. When $\mathrm{SSB}_{\rm now}>B_{\lim}>0$ and $F_{\rm now}>0$, define the dimensionless log-distance and its fishing-pressure rescaling by
\begin{equation}
R_B=\log\left(\frac{\mathrm{SSB}_{\rm now}}{B_{\lim}}\right),
\qquad
\Theta_F=\frac{R_B}{F_{\rm now}}.
\label{eq:fishing-only}
\end{equation}
The quantity $\Theta_F$, labelled the \textbf{Fishing-Only Time-to-Reference}, is the crossing time of the deliberately incomplete comparison process $\dot B=-F_{\rm now}B$. It is a removals-only pressure time scale: the time unit comes from rescaling a stock-reference margin by one isolated gross-loss rate. Because recruitment, somatic growth, maturation, natural mortality, density dependence, environmental forcing, and future policy are omitted, $\Theta_F$ is not a net biomass depletion diagnostic, a demographic hitting-time estimate, or a member of the $J_A^{\rm gross}$--$H_A^{\rm loc}$--$T_A$ hierarchy.

When $B(t)>B_{\lim}$, a genuinely local biomass-decline ratio would instead use the net biological balance,
\begin{equation}
H_B^{\rm loc}(t)
=\frac{B(t)-B_{\lim}}{[-\dot B(t)]_+},
\label{eq:biomass-local-ratio}
\end{equation}
with the same extended-real convention as Eq.~\eqref{eq:local-horizon}. Estimating it requires a compatible estimate of net $\dot B$, and a demographic hitting time requires a fully specified population model. A demographic application therefore requires at least adult and juvenile or recruit states, life-history transitions, natural mortality, and fishing selectivity. A support-pool application requires a measured abiotic donor stock and its balance. RAM Legacy SSB and $F$ data \cite{ricard2012} do not by themselves supply these quantities or models.

\section{From accounting to sustainability}

The applications above illustrate accounting diagnostics, not trajectory-level sustainability tests. The contemporaneous condition $b_i\ge0$ for every component is a service-adequacy condition. It is neither necessary nor sufficient for indefinite safety without further assumptions.

The ledger-to-viability reduction requires additional declared objects. For the six-compartment illustration, one may define a physical safe set
\[
K_x=\{z\in\R_+^6:z_q\ge z_q^{\min}\text{ for every thresholded coordinate }q\}
\]
and, for a demand path, the state-dependent operating set
\[
\mathcal U_d(z,t)
=\{u\in\mathcal U(z,t):\mathcal S_{\rm svc}(z,u,\theta)-d(t)\in\R_+^n\}.
\]
A policy-level question then asks whether an admissible $u=\pi(\text{available information})$ can remain in $\mathcal U_d$ while keeping the physical trajectory in $K_x$. Time-varying demand may instead be included in an augmented state and safe set. This construction links the incidence ledger, the service readout, and a viability problem, but it does not identify their certification statements.

For an actual policy $\pi$ and disturbance path $d$, trajectory safety means
\[
x^{\pi,d}(t;x_0)\in K\qquad\text{for all }t\ge0.
\]
Viability asks whether there exists at least one admissible policy with that property \cite{aubin1991,cury2005}; robust viability asks whether one causal policy works across a declared disturbance class. These concepts are part of the wider theory programme, not results computed in this accounting paper.

Stability must also remain separate from safety. A stable equilibrium may lie outside $K$, while a bounded non-equilibrium trajectory may remain inside it. Likewise, a long local depletion ratio does not certify viability, and a short gross turnover ratio does not prove net depletion.

\section{Limitations and research boundary}

\begin{enumerate}
    \item The matrix scaffold permits several conservation rows, but the worked equations instantiate one conserved material. A coupled multi-element application must specify its conservation matrix, stoichiometric conversions, and double-counting rules.
    \item The six compartments, two constant routing fractions, and absorbing-sink convention are illustrative constitutive choices. Feedback from accumulated sink mass and material remobilization require explicit rate dependencies or additional ledger columns.
    \item Service readouts, feasible demand domains, and thresholds are application-specific and may be uncertain, nonlinear, spatially distributed, or set-valued.
    \item The scalar counterexample is an elementary observation on an unrestricted balance space. Feasible-domain restrictions can alter what a scalar statistic implies and require a separate implication proof.
    \item Local ratios inherit measurement and derivative-estimation error and are not substitutes for trajectory simulation. The groundwater anomaly index and fisheries pressure time scale are classified outside physical net-stock forecasts.
    \item The applications do not constitute a calibrated joint demonstration of the full chain from typed ledger through service balances to viability. Computing all three diagnostic levels on one multi-compartment trajectory ensemble remains a useful next application.
    \item The paper does not model governance delays, endogenous effort dynamics, substitution technologies, spatial transfers, justice constraints, or exergy. Those modules should be added only when required by the claim under study.
\end{enumerate}

\section{Conclusion}

Componentwise accounting can support stronger sustainability claims only when its physical and logical roles are kept separate. A vector balance prevents compensation from being hidden by presentation, but a vector alone is not a general sustainability theory. Current service adequacy may remain support-dependent when the same service is infeasible under declared regenerative, boundary, quality, or exergy restrictions, yet that distinction alone does not establish net drawdown. A conservation identity does not guarantee physical non-negativity unless transfers are donor limited. Gross uptake demonstrates dependency, not decline. A local stock-to-rate ratio is not a forecast. Spawning biomass is not an abiotic support pool.

These distinctions yield a bounded accounting framework. The incidence matrix supplies typed physical balances, the service readout supplies contemporaneous balance vectors, and declared threshold and policy sets supply the inputs to a separate viability problem. Establishing a broader theory still requires application-specific observation, governance, uncertainty, and reduction structures; the six-compartment illustration is not itself that theory.

\section*{Data and code availability}
The groundwater calculations use the cited G3P product and the phosphate calculation uses the cited USGS reserve and production quantities. The accompanying computational supplement provides the extraction procedures, source-product versions, spatial masks, processed tables, and numerical settings used for these results.

\input{../shared/references.tex}

\end{document}
